\section{Introduction}
\label{sec:intro}

Nothing guarantees that a graph a large language model (LLM) pipeline builds from text can answer the questions people will ask. Ontology engineering approaches this with competency questions~\cite{gruninger1995methodology,noy2001ontology101}. Here a competency question (CQ) is a test question generated from a text chunk, with an expected answer span. Recent pipelines generate CQs from text, grow a schema from them and repair what fails. We call this cycle the repair loop: build the graph, write test questions, find the ones it cannot answer, and patch the graph, repeatedly (Section~\ref{sec:pipeline}).

Consider the question ``What is the cause of death of director of film Wedding Bills?'' from 2WikiMultihopQA (2Wiki)~\cite{ho2020twowiki}, multi-hop questions over Wikipedia. The source text says the director died ``from Parkinson's disease''. With Qwen, an open-weight LLM family, the graph built by the CQ-guided single pass (one extraction pass whose schema comes from 40 generated CQs) with no relation list has no cause-of-death triple for him, and the answerer says ``unknown''. Given a relation list, here 22 relation names that include ``cause of death'', the same pipeline extracts the triple and answers correctly. The relation list is a flat vocabulary of names, with no types or ranges, and is the variable this paper isolates.

Published CQ-driven pipelines~\cite{feng2024ontologygrounded,kaiser2025hydra,choi2026socratickg} and loops~\cite{huang2025tcrqf,chong2026evoontology,zhang2026oak} change several components at once and report one final score, which cannot say where a gain comes from. We rebuild their components in one pipeline family rather than rerun each system. This study scores finished graphs on held-out questions and tests the repair step against random patches of the same size.

We build a graph from text, then test it with questions it never saw. A fact helps held-out question answering (QA) only if the graph contains it and retrieval returns it: removing the gold-evidence facts from what the answerer sees lowers answer correctness by 40 points or more (Appendix~\ref{app:reach}). Our main measure, grounded QA, is the share of held-out questions answered correctly from retrieved graph facts that contain the answer, scored by code; answer correctness drops the support requirement, and the SPARQL pass rate measures querying the graph. The answerer sees at most 60 retrieved facts (the answer fact cap). The corpora are 2Wiki, a T-REx~\cite{elsahar2018trex} movie subset with one-hop questions, and HotpotQA~\cite{yang2018hotpotqa}. Values are means over shards (groups of questions with their source paragraphs; a unit is one shard built under one run seed), and differences are in points with 95\% intervals over shards; two arms are \emph{level} when that interval includes zero.

An arm is a pipeline variant, and the relation list is the variable that separates them. We compare these arms (Table~\ref{tab:arms} gives the full definitions):
\begin{itemize}[nosep]
  \item \emph{CQ-guided single pass}: one extraction pass whose schema comes from 40 generated CQs.
  \item \emph{Repair loop}: the CQ-guided single pass plus rounds of new CQs that test and patch the graph.
  \item \emph{Plain extraction}: every chunk extracted with free-form predicates, no CQs, no list.
  \item \emph{Seeded extraction}: plain extraction with the relation list in every prompt; predicates stay free-form.
  \item \emph{List-schema pass}: one pass whose schema is only the 22 list names, no CQs; predicates are restricted to those names.
  \item \emph{Combined pipeline} (our full list-plus-CQ pipeline): the list, CQs for every chunk, extra extraction passes, CQ evaluation.
\end{itemize}
Our main result is that on 2Wiki the relation list, not the chunk-generated CQs or the repair loop, produces the held-out QA gain. Three findings follow.

(1) On 2Wiki, a list-schema pass matches the combined pipeline built with the same list in answer correctness at one eighth of its construction tokens, and leads it in grounded QA when the answerer sees up to 60 facts: by 4.6 points with the 22 most frequent gold relation types and by 5.3 with a list induced without gold (60 units each). At 120 facts, where no cap binds, the two are level. A single pass whose schema comes from generated CQs, given the same list, scores 11.1 points lower; CQs pay off on the questions their own generator writes, where the combined pipeline leads by 18.1 points.

(2) Seeded extraction costs 6.7 to 8.3 times fewer tokens than the repair loop on 2Wiki, 7.1 on HotpotQA and 16 to 17 on T-REx, but it stores its facts under hundreds of relation names and passes only 7.5\% of the SPARQL tests; mapping its predicates to the 22 list names recovers about half of the lost SPARQL passes.

(3) The repair loop's lead depends on question type and corpus: pooled over 60 2Wiki shards it leads seeded extraction by 3.2 points in grounded QA, all of it from bridge-comparison questions (it trails on comparison questions); on HotpotQA seeded extraction leads at every fact cap, and an induced list lowers plain extraction by 4.0 points ($+$0.8 on 2Wiki).

We release plans, deviation logs, code and graphs.\footnote{\releaseurl}
